\uuid{NXTe}
\titre{Raisonnements}
\auteur{bodin ; croizat ; sacre}
\organisation{exo7}
\licence{CC BY-NC-SA 4.0}
\datecreate{2022-03-08}
\niveau{L1}
\module{Algèbre}
\chapitre{Logique, ensemble, raisonnement}
\sousChapitre{Absurde et contraposée}
\difficulte{1}
\theme{Raisonnements}
\exo7qcmid{hilisit.logique.041}

\contenu{
\begin{question}
\qclassification{100.03}
Pour montrer que pour tout $x\in\Rr$ et tout $n\in\Nn$ on a $(1+x)^n \ge 1+nx$, quelle est la démarche la plus adaptée ?
\begin{answers}
	\good{On fixe $x$, on fait une récurrence sur $n$.}
	\bad{On fixe $n$, on fait une récurrence sur $x$.}
	\bad{Par l'absurde on suppose $(1+x)^n < 1+nx$.}
	\bad{Par disjonction des cas $n$ pair/$n$ impair.}
\end{answers}
\begin{explanations}
On fixe $x \in \Rr$, on fait une récurrence sur $n \in \Nn$. On ne peut pas procéder à une récurrence sur $x \in \Rr$.
\end{explanations}
\end{question}
}
